\scanpage{239}
\mspage{189}

On utilise la \deleted{fibration} $W_\infty$-fibration (parfaite)

\diagram{\begin{tikzcd}[ampersand replacement=\&,column sep=0.25em,row sep=1.2em]
E\arrow[d] \& {}=\underline{\mathrm{Ch}}(X,x,{-}) \& \qquad E\ \text{$W_\infty$-asphérique}\\
X \& \&
\end{tikzcd}}

de fibre \inserted{$Z$} en $x$ égale à $\Omega(X,x)$.

Supposons $a_n)$ pour $X$, donc $a_{n-1})$ pour $\Omega(X,x)$. Par hyp. de récurrence, on a donc

\[
H^{\unsure{i|n}}(Z,\deleted{\unsure{A}}M)=0\qquad\text{pour }\deleted{\unsure{q}}\,i\in[1,n-1].
\]

$\deleted{\unsure{A}}M$ \deleted{faisceau constant sur $Z$} groupe abélien. Donc dans la suite spectrale de Leray

\[
H^{\unsure{*|n}}(E,M)\Longleftarrow E^{pq}_2=H^p(X,H^q(Z,M))
\]

on a

\[
E^{pq}_2=0\quad\text{si}\quad q\in[1,n-1]
\]

\diagram{\begin{tikzcd}[ampersand replacement=\&,column sep=.36em,row sep=.24em]
q \& \& \& \& \&\\
|[label=left:{$\scriptstyle n$}]| \bullet\arrow[u] \& \bullet \& \bullet \& \bullet \& \bullet \& \cdots\\
{} \& 0 \& \& 0 \& \&\\
\bullet\arrow[uu,no head]\arrow[rrrrr] \& \bullet \& \bullet \& \bullet \& \bullet \& {} \\
\& \& \& \& \& p

\end{tikzcd}}

On a donc la suite exacte

\[
0\longrightarrow H^n(X,M)\longrightarrow H^n(E,M)
\]

et comme $E$ est $W_\infty$-asphérique, on a $H^n(E,M)=0$, donc $H^n(X,M)=0$.

Inversement, supposons $b_n)$, prouvons $a_n$. On sait déjà qu'on a $a_{n-1}$, donc $a_{n-2}$ pour $\Omega(X,x)$, \deleted{\unsure{i.e.}} qui implique

\[
H^i(Z,M)=0\qquad\text{pour}\qquad i\in[1,n-2],
\]

\[
\text{donc}\qquad E^{pq}_2=0\quad\text{si}\quad q\in[1,n-2].
\]

On a donc la suite exacte
% archive sheet number: 239
